\SourceLeaf{336}

\section*{[II.]}
\begin{center}VORBEREITUNGEN ZUR ALLGEMEINEN THEORIE\\DER BIQUADRATISCHEN RESTE.\end{center}
\pdfbookmark[1]{II. Vorbereitungen zur allgemeinen Theorie der biquadratischen Reste}{nachlass-biquadratic-II}
\BookRule{.15\textwidth}
\par\medskip\pdfbookmark[2]{(1.)}{nachlass-biquadratic-II-prefatory-1}\centerline{(1.)}\smallskip
Es sei $P=x+iy$, wo weder $x$ noch $y$ eine ganze Zahl ist. Wir bezeichnen die. Zahl $+1$ durch $LP,L'P,L''P,L'''P$, je nachdem $P$ im ersten, zweiten, dritten oder vierten Quadranten liegt (im ersten und zweiten Quadranten ist $[y]$ gerade, im dritten und vierten ungerade; im ersten und vierten ist $[x]$ gerade, im zweiten und dritten ungerade). In allen Fällen, wo diese Zeichen nicht $=1$ sind, werden sie $=0$ vorausgesetzt. Man hat dann folgende $24$ Relationen

\[
\begin{array}{lll}
L(P\pm1)=L'P&L(P\pm i)=L'''P&L(P\pm1\pm i)=L''P\\
L'(P\pm1)=LP&L'(P\pm i)=L''P&L'(P\pm1\pm i)=L'''P\\
L''(P\pm1)=L'''P&L''(P\pm i)=L'P&L''(P\pm1\pm i)=LP\\
L'''(P\pm1)=L''P&L'''(P\pm i)=LP&L'''(P\pm1\pm i)=L'P
\end{array}
\]
\[
\begin{array}{lll}
LiP=L'''P&L(-P)=L''P&L(-iP)=L'P\\
L'iP=LP&L'(-P)=L'''P&L'(-iP)=L''P\\
L''iP=L'P&L''(-P)=LP&L''(-iP)=L'''P\\
L'''iP=L''P&L'''(-P)=L'P&L'''(-iP)=LP
\end{array}
\]

\par\medskip\pdfbookmark[2]{(2.)}{nachlass-biquadratic-II-prefatory-2}\centerline{(2.)}\smallskip
Durch $PP'$ oder $z$ bezeichnen wir eine Linie, die von $P$ anfängt und in $P'$ endigt. Sie braucht nicht gerade zu sein. Wir legen allen geraden Linien von $2x+2iy$ nach $2x+(2y+1)i$ gezogen (wo $x,y$ indefinite alle ganzen Zahlen bedeuten) eine positive und eine negative Seite bei; für jene wählen wir die rechte, für diese die linke. Durch $Tz$ bezeichnen wir die Anzahl aller Schnitte
\SourceLeaf{337}

der Linie $z$ mit den eben gedachten Linien, als positiv gezählt diejenigen, wo $z$ von der negativen Seite auf die positive übergeht, als negativ die andern. Ferner setzen wir
\[
Tz-T(z-1)=Sz
\]

($z-1$ ist eine der $z$ parallele Linie, die von dem Punkte $P-1$ nach $P'-1$ geht). Offenbar brauchen wir nur dem oben gedachten System von Linien noch die von $2x+1+2yi$ nach $2x+1+(2y+1)i$ gezogenen beizufügen und deren linke Seiten positiv und die rechten als negativ zu betrachten, um in $Sz$ die Anzahl aller Schnitte von $z$ mit diesem zweifachen System von Geraden zu erkennen. Wir haben nun ferner
\[
\begin{aligned}
T(-z)&=-T(z+i)\\
T(z)+T(z+i)&=[\tfrac12x]-[\tfrac12x']\\
S(z+1)&=-Sz\\
S(z+i)&=-Sz+LP+L'''P-LP'-L'''P'\\
S(z+1+i)&=Sz-LP+L'''P+LP'+L'''P'\\
Siz&=Sz-LP+LP'\\
S(-z)&=Sz-LP-L'''P+LP'+L'''P'\\
S(-iz)&=Sz+L'P-L'P'
\end{aligned}
\]

\renewcommand{\Article}[1]{\par\medskip\Needspace{4\baselineskip}\pdfbookmark[2]{Art. #1}{nachlass-biquadratic-II-art-#1}\centerline{\textbf{#1.}}\nobreak\smallskip}
\Article{1}
Wir betrachten in der Ebene zwei Gattungen von Punkten; einmal die, denen ganze Zahlen entsprechen; dann diejenigen, welche durch Producte aus ganzen Zahlen in die Grösse $Q=\frac{m}{2M}$ bestimmt werden. Wir können dieselben durch die Benennungen Punkte der ersten und Punkte der zweiten Ordnung unterscheiden.

\Article{2}
Indem wir jeden Punkt der zweiten Ordnung mit seinen vier Nachbarn durch gerade Linien verbinden, die wir \emph{Ligaturen} nennen werden, theilt sich die ganze Ebene in unendlich viele Quadrate. Die Punkte der ersten Ordnung liegen theils innerhalb dieser Quadrate, theils auf den Ligaturen innerhalb der 
\SourceLeaf{338}

Grenzen derselben, theils auf den Grenzen der Ligaturen, das letzte, wenn sie zugleich Punkte der zweiten Ordnung sind. Ist $kQ$ ein solcher Punkt, so muss insofern $m,M$ ohne gemeinschaftlichen Theiler und beide ungerade sind, $k$ durch $M$ theilbar sein.

\Article{3}
Bei den Ligaturen können wir zugleich einen Unterschied zwischen dem Anfangspunkte und Endpunkte machen, also $PQ$ von $QP$ unterscheiden, oder auch in einigen Fällen diesen Unterschied bei Seite setzen. Wir nennen zwei solche Ligaturen entgegengesetzte. Bezeichnen können wir überhaupt am bequemsten die Ligaturen durch ihren Anfangs- und Endpunkt, die man allenfalls in eine Klammer einschliessen mag. Einer Ligatur entgegengesetzte soll durch das doppelte Ueberstreichen angedeutet werden $QP=\overline{\overline{PQ}}$.

\Article{4}
Jedes der gedachten Quadrate wird von vier solchen Ligaturen einge\-schlossen
\[
\begin{gathered}
\{kQ,(k+1)Q\},\ \{(k+1)Q,(k+1+i)Q\},\\
\{(k+1+i)Q,(k+i)Q\},\ \{(k+i)Q,kQ\}\dotsfour\Omega
\end{gathered}
\]
denen es zur rechten liegt. Es ist wichtig hiebei auf die Form der Zahl $k$ zu sehen, und wir unterscheiden in dieser Beziehung viererlei Quadrate, je nachdem $k\equiv0,1,1+i,i\ (\mathrm{mod.}\ 2)$ ist, und bedienen uns dann der Zahlen $0,1,2,3$, die wir resp. die Intensoren der Quadrate nennen.

\Article{5}
Den Ligaturen legen wir dieselben Intensoren bei, welche die ihnen zur rechten liegenden Quadrate haben.

\Article{6}
Wir haben nun ein anderes grösseres Quadrat $\Omega'$ zu betrachten, nemlich dasjenige. welches entsteht, wenn das in $4$ angezeigte für $k=0$, mit $M$ multiplicirt wird: dies wird also durch die geraden Linien $\mu,\mu',\mu'',\mu'''$ begrenzt
\[
\{0,\tfrac12m\},\ \{\tfrac12m,\tfrac12(1+i)m\},\ \{\tfrac12(1+i)m,\tfrac12im\},\ \{\tfrac12im,0\}
\]
Es besteht aus ganzen Quadraten $\Omega$ und Stücken solcher Quadrate; man zähle
\SourceLeaf{339}

alle Punkte der ersten Ordnung innerhalb desselben zusammen, indem man für jeden Punkt den Intensor des Quadrats $\Omega$, worin er liegt, nimmt, diese Summe oder deren kleinster Rest nach dem Modulus $4$ heisst der Decident von $M$ für den Modulus $m$, und bestimmt die biquadratische Modalität von $M$ in Beziehung auf diesen Modulus.

\Article{7}
Wir zerlegen das Quadrat $\Omega'$ in $5$ Stücke auf folgende Art. Man verbinde den Punkt $0$ mit $\frac14(1+i)(m-1)$ durch die Linie $\lambda$, die durch lauter Ligaturen innerhalb $\Omega'$ gehe. Es sei
\[
\tfrac12m+i\lambda=\lambda',\quad\tfrac12(1+i)m-\lambda=\lambda'',\quad\tfrac12im-i\lambda=\lambda'''
\]
diese $4$ Linien gehen also von den Ecken des Quadrats $\Omega'$ aus ins Innere und endigen sich an den vier Ecken des innersten Quadrats, dessen Intensor $0$ sein wird, wenn $m\equiv1\ (\mathrm{mod.}\ 2+2i)$; die Ligaturen dieses Quadrats seien $\nu,\nu',\nu'',\nu'''$.

Die $5$ Stücke werden also begrenzt sein
\[
\begin{array}{rll}
\mathrm{I}.&\ldots&\mu,\ \lambda',\ \overline{\overline\nu},\ \overline{\overline\lambda}\\
\mathrm{II}.&\ldots&\mu',\ \lambda'',\ \overline{\overline{\nu'}},\ \overline{\overline{\lambda'}}\\
\mathrm{III}.&\ldots&\mu'',\ \lambda''',\ \overline{\overline{\nu''}},\ \overline{\overline{\lambda''}}\\
\mathrm{IV}.&\ldots&\mu''',\ \lambda,\ \overline{\overline{\nu'''}},\ \overline{\overline{\lambda'''}}\\
\mathrm{V}.&\multicolumn{2}{l}{\text{das innere Quadrat }\nu,\ \nu',\ \nu'',\ \nu'''}
\end{array}
\]
Der Decident ist also die Aufzählung aller Punkte erster Ordnung in I. II. III. IV.

\Article{8}
Der Kürze wegen soll Intensor irgend eines Punkts der Intensor des Quadrats sein, in dem er liegt, und durch vorgesetztes $\Upsilon$ ausgedrückt werden.

\Article{9}
Der Decident ist also
\[
\Sigma\Upsilon P+\Sigma\Upsilon P'+\Sigma\Upsilon P''+\Sigma\Upsilon P'''
\]
wo $P$ alle Punkte in I. u. s. w. bedeuten.

\SourceLeaf{340}

\Article{10}
Wir betrachten nun noch den Raum $\mathrm{VI}=-i\mathrm{IV}$, welcher ausserhalb $\Omega'$ liegt, sich aber durch $\mu$ an I anschliesst und mit ihm zusammen den Raum $\omega$ ausmacht, der aus $AA+BB$ vollständigen Quadraten besteht. Bedeutet $\Pi$ alle ganzen; $\Pi'$ alle um $\frac12i$ vermehrten ganzen Punkte dieses Raumes, so lässt sich leicht beweisen, dass der Decident
\[
\begin{aligned}
={}&\Sigma\Upsilon\Pi-\Sigma\Upsilon\Pi'+\text{Anzahl aller ganzen Punkte innerhalb VI}\\
&-\text{Anzahl aller halben Punkte innerhalb VI.}
\end{aligned}
\]
